%% This is file `elsarticle-template-1-num.tex',
%%
%% Copyright 2009 Elsevier Ltd
%%
%% This file is part of the 'Elsarticle Bundle'.
%% ---------------------------------------------

\documentclass[a4paper]{article}

\usepackage{graphicx}
\usepackage{import}
\usepackage{xifthen}
\usepackage{pdfpages}
\usepackage{transparent}
\usepackage{cancel}
\usepackage{amssymb}
\usepackage{amsmath}
\usepackage{amsthm}
\usepackage{ mathrsfs }
\usepackage{ dsfont }
\usepackage{tikz-cd}
\usepackage[a4paper,top=3cm,bottom=2cm,left=3cm,right=3cm,marginparwidth=1.75cm]{geometry}
\usepackage{float}
\usepackage{lineno}
\usepackage{hyperref}
\usepackage{enumitem}

\newtheorem{name}{Printed output}
\newtheorem{mydef}{Definition}
\newtheoremstyle{break}% name
  {}%
  {}%
  {}%
  {}%
  {\bfseries}%
  {.}%
  {\newline}%
  {}

\theoremstyle{plain}
\newtheorem{question}{Question}

\newtheorem{theorem}{Theorem}[section]
\newtheorem*{lemma*}{Lemma}
\newtheorem{lemma}[theorem]{Lemma}
\newtheorem{corollary}[theorem]{Corollary}
\newtheorem{proposition}[theorem]{Proposition}

\theoremstyle{definition}
\newtheorem{definition}[theorem]{Definition}
\newtheorem*{note*}{Note}
\newtheorem*{claim}{Claim}
\newtheorem*{remark}{Remark}
\newtheorem{example}{Example}
\newtheorem*{solution}{Solution}
\newtheorem{exercise}{Exercise}

\numberwithin{equation}{subsection}

\newcommand{\incfig}[2][1]{%
    \def\svgwidth{#1\columnwidth}%
    \import{./Figures/}{#2.pdf_tex}%
}
\newcommand{\hodge}{{*}}
\newcommand{\tang}{\operatorname{\textbf{\textup{t}}}}
\newcommand{\norm}{\operatorname{\textbf{\textup{n}}}}
\newcommand{\R}{\mathbb{R}}
\newcommand{\bP}{\mathbb{bP}}
\newcommand{\N}{\mathbb{N}}
\newcommand{\Q}{\mathbb{Q}}
\newcommand{\A}{\mathbb{A}}
\newcommand{\bS}{\mathbb{S}}
\newcommand{\Z}{\mathbb{Z}}
\newcommand{\F}{\mathbb{F}}
\newcommand{\B}{\mathbb{B}}
\newcommand{\T}{\mathbb{T}}
\newcommand{\K}{\mathbb{K}}
\newcommand{\C}{\mathbb{C}}
\newcommand{\D}{\mathbb{D}}

\usepackage{comment}
\newcommand{\bmv}{\bm{v}}
\newcommand{\bmu}{\bm{u}}
\newcommand{\bmw}{\bm{w}}
\newcommand{\bmx}{\bm{x}}

\newcommand{\cF}{\mathcal{F}}
\newcommand{\cA}{\mathcal{A}}
\newcommand{\cM}{\mathcal{M}}
\newcommand{\cU}{\mathcal{U}}
\newcommand{\cB}{\mathcal{B}}
\newcommand{\cD}{\mathcal{D}}
\newcommand{\cL}{\mathcal{L}}
\newcommand{\cS}{\mathcal{S}}
\newcommand{\cG}{\mathcal{G}}
\newcommand{\cT}{\mathcal{T}}
\newcommand{\cC}{\mathcal{C}}
\newcommand{\cH}{\mathcal{H}}
\newcommand{\cO}{\mathcal{O}}
\newcommand{\cJ}{\mathcal{J}}

\newcommand{\frp}{\mathfrak{p}}
\newcommand{\frq}{\mathfrak{q}}
\newcommand{\frm}{\mathfrak{m}}
\newcommand{\frP}{\mathfrak{P}}

\newcommand{\msF}{\mathscr{F}}
\newcommand{\msA}{\mathscr{A}}
\newcommand{\msM}{\mathscr{M}}
\newcommand{\msE}{\mathscr{E}}
\newcommand{\msU}{\mathscr{U}}
\newcommand{\msB}{\mathscr{B}}
\newcommand{\msD}{\mathscr{D}}
\newcommand{\msL}{\mathscr{L}}
\newcommand{\msS}{\mathscr{S}}
\newcommand{\msG}{\mathscr{G}}
\newcommand{\msT}{\mathscr{T}}
\newcommand{\msC}{\mathscr{C}}
\newcommand{\msH}{\mathscr{H}}
\newcommand{\msO}{\mathscr{O}}
\newcommand{\msJ}{\mathscr{J}}
\newcommand{\msP}{\mathscr{P}}

\newcommand{\vect}[1]{\underline{\textbf{#1}}}
\newcommand{\llangle}{\left\langle}
\newcommand{\rrangle}{\right\rangle}
\newcommand{\llbrace}{\left\lbrace}
\newcommand{\rrbrace}{\right\rbrace}
\newcommand{\setof}[1]{\left\lbrace #1 \right\rbrace}

\newcommand{\cFA}{\cF(\cA)}
\newcommand{\sigmaA}{\sigma(\cA)}
\newcommand{\del}[2]{\frac{\partial #1}{\partial #2}}

\newcommand{\innerproductb}[2]{\llangle\bm{#1},\bm{#2}\rrangle}
\newcommand{\innerproduct}[2]{\llangle{#1},{#2}\rrangle}

\newcommand{\cl}{\operatorname{cl}}
\newcommand{\coker}{\operatorname{coker}}
\newcommand{\re}{\operatorname{Re}}
\newcommand{\im}{\operatorname{Im}}
\newcommand{\dist}{\operatorname{dist}}
\newcommand{\Span}{\operatorname{Span}}
\newcommand{\Supp}{\operatorname{Supp}}
\newcommand{\id}{\operatorname{id}}
\newcommand{\sh}{\operatorname{sh}}
\newcommand{\var}{\operatorname{Var}}
\newcommand{\Proj}{\operatorname{Proj}}
\newcommand{\Ker}{\operatorname{Ker}}
\newcommand{\Mor}{\operatorname{Mor}}
\newcommand{\supp}{\operatorname{supp}}
\newcommand{\Ran}{\operatorname{Ran}}
\newcommand{\Res}{\operatorname{Res}}
\newcommand{\Mod}{\operatorname{Mod}}
\newcommand{\res}{\operatorname{res}}
\newcommand{\Frac}{\operatorname{Frac}}
\newcommand{\End}{\operatorname{End}}
\newcommand{\Sym}{\operatorname{Sym}}
\newcommand{\Spec}{\operatorname{Spec}}
\newcommand*{\sheafhom}{\mathcal{H}\kern -.5pt \it{ om}}
\newcommand*{\sheafspec}{\mathcal{S}\kern -.5pt \it{ pec}}
\newcommand*{\sheafproj}{{\mathcal{P}\kern -.5pt \it{ roj}}}
\newcommand{\Div}{\operatorname{div}}
\newcommand{\Aut}{\operatorname{Aut}}
\newcommand{\Hom}{\operatorname{Hom}}
\newcommand{\Qcoh}{\operatorname{Qcoh}}
\newcommand{\PGL}{\operatorname{PGL}}
\newcommand{\Op}{\operatorname{Op}}
\newcommand{\WF}{\operatorname{WF}}
\newcommand{\Char}{\operatorname{Char}}
\newcommand{\Ell}{\operatorname{Ell}}
\newcommand{\RE}{\operatorname{Re}}
\newcommand{\IM}{\operatorname{Im}}
\newcommand{\cosec}{\operatorname{cosec}}

\setlength{\parindent}{0pt}
\setlength{\parskip}{1em}

\begin{document}

\title{\textbf{Problem Sheet: Chain Rule and Substitutions 2 Solution}}
\author{\textbf{Jeffrey Thompson}}
\maketitle

\begin{exercise}
Differentiate the following functions.
\begin{enumerate}
    \item $\sin(\cos(3x))$
    \item $e^{2\sin(4x)}$
    \item $\cos(\ln(5x))$
    \item $(\ln(2x))^3$
\end{enumerate}
\end{exercise}

\begin{solution}

\begin{enumerate}
    \item Differentiating $\sin(\cos(3x))$

    Let
    \[
        y = \sin(\cos(3x)), \qquad u = \cos(3x)
    \]
    so that $y = \sin u$.

    Then
    \begin{align*}
        \frac{dy}{du} &= \cos u \\
        \frac{du}{dx} &= -3\sin(3x)
    \end{align*}

    Therefore
    \begin{align*}
        \frac{dy}{dx} &= \frac{dy}{du}\frac{du}{dx}\\
        &= \cos u \cdot (-3\sin(3x))\\
        &= -3\sin(3x)\cos(\cos(3x))
    \end{align*}

    \item Differentiating $e^{2\sin(4x)}$

    Let
    \[
        y = e^{2\sin(4x)}, \qquad u = 2\sin(4x)
    \]
    so that $y = e^u$.

    Then
    \begin{align*}
        \frac{dy}{du} &= e^u \\
        \frac{du}{dx} &= 8\cos(4x)
    \end{align*}

    Therefore
    \begin{align*}
        \frac{dy}{dx} &= \frac{dy}{du}\frac{du}{dx}\\
        &= e^u \cdot 8\cos(4x)\\
        &= 8e^{2\sin(4x)}\cos(4x)
    \end{align*}

    \item Differentiating $\cos(\ln(5x))$

    Let
    \[
        y = \cos(\ln(5x)), \qquad u = \ln(5x)
    \]
    so that $y = \cos u$.

    Then
    \begin{align*}
        \frac{dy}{du} &= -\sin u \\
        \frac{du}{dx} &= \frac{1}{x}
    \end{align*}

    Therefore
    \begin{align*}
        \frac{dy}{dx} &= \frac{dy}{du}\frac{du}{dx}\\
        &= -\sin u \cdot \frac{1}{x}\\
        &= -\frac{\sin(\ln(5x))}{x}
    \end{align*}

    \item Differentiating $(\ln(2x))^3$

    Let
    \[
        y = (\ln(2x))^3, \qquad u = \ln(2x)
    \]
    so that $y = u^3$.

    Then
    \begin{align*}
        \frac{dy}{du} &= 3u^2 \\
        \frac{du}{dx} &= \frac{1}{x}
    \end{align*}

    Therefore
    \begin{align*}
        \frac{dy}{dx} &= \frac{dy}{du}\frac{du}{dx}\\
        &= 3u^2 \cdot \frac{1}{x}\\
        &= \frac{3(\ln(2x))^2}{x}
    \end{align*}
\end{enumerate}
\qed
\end{solution}

% ------------------------------------------------------------

\begin{exercise}
Let $f(x)=e^{3\sin x} + \sin(x^2)$. Find $f'(x)$ and $f''(x)$.
\end{exercise}

\begin{solution}

Let
\[
    f(x) = e^{3\sin x}, \qquad g(x) = \sin(x^2).
\]
Then
\[
    \left(e^{3\sin x} + \sin(x^2)\right)' = f'(x) + g'(x).
\]

For $f(x)$, let $y = e^{3\sin x}$ and $u = 3\sin x$. Then $y = e^u$, so
\begin{align*}
    \frac{dy}{du} &= e^u\\
    \frac{du}{dx} &= 3\cos x
\end{align*}
and therefore
\[
    f'(x) = 3e^{3\sin x}\cos x.
\]

For $g(x)$, let $y = \sin(x^2)$ and $u = x^2$. Then $y = \sin u$, so
\begin{align*}
    \frac{dy}{du} &= \cos u\\
    \frac{du}{dx} &= 2x
\end{align*}
and therefore
\[
    g'(x) = 2x\cos(x^2).
\]

Hence the first derivative is
\[
    \left(e^{3\sin x} + \sin(x^2)\right)' = 3e^{3\sin x}\cos x + 2x\cos(x^2).
\]

Now differentiate again.

For $3e^{3\sin x}\cos x$, use the product rule. Let
\[
    u = 3e^{3\sin x}, \qquad v = \cos x.
\]
Then
\[
    v' = -\sin x
\]
and to find $u'$ we again use the chain rule:
\[
    u' = 3\left(e^{3\sin x}\right)' = 3\left(3e^{3\sin x}\cos x\right) = 9e^{3\sin x}\cos x.
\]
Therefore
\begin{align*}
    (uv)' &= uv' + u'v\\
    &= 3e^{3\sin x}(-\sin x) + 9e^{3\sin x}\cos x \cdot \cos x\\
    &= -3e^{3\sin x}\sin x + 9e^{3\sin x}\cos^2 x
\end{align*}

For $2x\cos(x^2)$, use the product rule again. Let
\[
    u = 2x, \qquad v = \cos(x^2).
\]
Then
\[
    u' = 2
\]
and for $v'$ let $w = x^2$, so that $v = \cos w$. Then
\begin{align*}
    \frac{dv}{dw} &= -\sin w\\
    \frac{dw}{dx} &= 2x
\end{align*}
so
\[
    v' = -2x\sin(x^2).
\]
Therefore
\begin{align*}
    (uv)' &= uv' + u'v\\
    &= 2x(-2x\sin(x^2)) + 2\cos(x^2)\\
    &= -4x^2\sin(x^2) + 2\cos(x^2)
\end{align*}

Hence the second derivative is
\begin{align*}
    \left(3e^{3\sin x}\cos x + 2x\cos(x^2)\right)'
    &= -3e^{3\sin x}\sin x + 9e^{3\sin x}\cos^2 x \\
    &\phantom{=} -4x^2\sin(x^2) + 2\cos(x^2)
\end{align*}

\qed
\end{solution}

% ------------------------------------------------------------

\begin{exercise}
For each integral below, state a suitable substitution before evaluating the indefinite integral.
\begin{enumerate}
    \item $\displaystyle \int x e^{x^2}\,dx$
    \item $\displaystyle \int \frac{4x}{1+2x^2}\,dx$
    \item $\displaystyle \int 6x\sqrt{1-3x^2}\,dx$
\end{enumerate}
\end{exercise}

\begin{solution}

\begin{enumerate}
    \item Integrating $\displaystyle \int x e^{x^2}\,dx$

    We use the substitution
    \[
        u = x^2.
    \]

    Then
    \[
        du = 2x\,dx
    \]
    so
    \[
        x\,dx = \frac{1}{2}du.
    \]

    Therefore
    \begin{align*}
        \int x e^{x^2}\,dx
        &= \frac{1}{2}\int e^u\,du\\
        &= \frac{1}{2}e^u + c\\
        &= \frac{1}{2}e^{x^2} + c.
    \end{align*}

    \item Integrating $\displaystyle \int \frac{4x}{1+2x^2}\,dx$

    We use the substitution
    \[
        u = 1+2x^2.
    \]
    Then
    \[
        du = 4x\,dx.
    \]

    Therefore
    \begin{align*}
        \int \frac{4x}{1+2x^2}\,dx
        &= \int \frac{1}{u}\,du\\
        &= \ln|u| + c\\
        &= \ln|1+2x^2| + c.
    \end{align*}

    \item Integrating $\displaystyle \int 6x\sqrt{1-3x^2}\,dx$

    We use the substitution
    \[
        u = 1-3x^2.
    \]
    Then
    \[
        du = -6x\,dx,
    \]
    so
    \[
        6x\,dx = -du.
    \]

    Therefore
    \begin{align*}
        \int 6x\sqrt{1-3x^2}\,dx
        &= -\int \sqrt{u}\,du\\
        &= -\int u^{1/2}\,du\\
        &= -\frac{2}{3}u^{3/2} + c\\
        &= -\frac{2}{3}(1-3x^2)^{3/2} + c.
    \end{align*}
\end{enumerate}
\qed
\end{solution}

\end{document}
